\chapter{Vectors and the Inner Product}

Economic data almost never comes as a single number: a household consumes a
bundle of goods, a firm produces several outputs, a portfolio holds many
assets. The natural object for such lists is the vector, and the single most useful operation on vectors is the inner
product, which turns two
vectors into one number and thereby measures expenditure, length, angle, and
orthogonality. This chapter drills vector arithmetic and the inner product;
it is the required review material for everything that follows. The other
half of the review topic (linear (in)dependence) has its own chapter
later in the book (Chapter~6); here we only prepare the ground with linear
combinations.

\section{Core Rules}

\begin{rulebox}{Vectors in $\R^{n}$}
A vector $x \in \R^{n}$ is an ordered list of $n$ real numbers, written as
a column:
\[
x = \begin{pmatrix} x_{1} \\ x_{2} \\ \vdots \\ x_{n} \end{pmatrix},
\qquad
x^{T} = (x_{1}, x_{2}, \dots, x_{n}) \text{ is the row version.}
\]
The zero vector $0$ has every entry $0$. The standard basis vectors
$e_{1}, \dots, e_{n}$ have a $1$ in one position and $0$ elsewhere. Two
vectors are equal iff they have the same dimension and all entries agree.
\end{rulebox}

\begin{rulebox}{Addition and scalar multiplication}
For $x, y \in \R^{n}$ and $\lambda \in \R$, both operations work entry by
entry:
\[
(x + y)_{i} = x_{i} + y_{i}, \qquad (\lambda x)_{i} = \lambda \, x_{i}.
\]
Addition requires equal dimensions. Geometric reading: a vector is an arrow
from the origin; $x + y$ is the tip-to-tail (parallelogram) construction;
$\lambda x$ stretches the arrow by $|\lambda|$ and, if $\lambda < 0$, flips
its direction. As with matrices, addition is commutative and associative,
and $\lambda(\mu x) = (\lambda \mu) x$.
\end{rulebox}

\begin{rulebox}{Linear and convex combinations}
A linear combination of $v_{1}, \dots, v_{k}$ is
any vector of the form
\[
\lambda_{1} v_{1} + \lambda_{2} v_{2} + \cdots + \lambda_{k} v_{k},
\qquad \lambda_{1}, \dots, \lambda_{k} \in \R.
\]
Every $x \in \R^{n}$ is the linear combination
$x = x_{1} e_{1} + \cdots + x_{n} e_{n}$ of the standard basis. If all
weights satisfy $\lambda_{i} \ge 0$ and $\lambda_{1} + \cdots + \lambda_{k}
= 1$, the combination is \emph{convex}, a
weighted average of the $v_{i}$.
\end{rulebox}

\begin{rulebox}{Inner product}
For $x, y \in \R^{n}$ the inner product is the scalar
\[
x \cdot y \;=\; x^{T} y \;=\; \sum_{i=1}^{n} x_{i} y_{i}
\;=\; x_{1} y_{1} + x_{2} y_{2} + \cdots + x_{n} y_{n}.
\]
Rules, for all $x, y, z \in \R^{n}$ and $\lambda \in \R$:
\begin{itemize}[leftmargin=1.6em,itemsep=0.3ex]
  \item Symmetry: $x \cdot y = y \cdot x$.
  \item Linearity in each slot: $(x + y) \cdot z = x \cdot z + y \cdot z$
        and $(\lambda x) \cdot y = \lambda \, (x \cdot y)$.
  \item Positive definiteness: $x \cdot x \ge 0$, with $x \cdot x = 0$ iff
        $x = 0$.
\end{itemize}
\end{rulebox}

\begin{rulebox}{Norm and unit vectors}
The (Euclidean) norm of $x \in \R^{n}$ is its length,
\[
\|x\| \;=\; \sqrt{x \cdot x} \;=\; \sqrt{x_{1}^{2} + \cdots + x_{n}^{2}},
\qquad \text{so } \|x\|^{2} = x \cdot x = x^{T} x .
\]
Properties: $\|x\| \ge 0$ with equality iff $x = 0$;
$\|\lambda x\| = |\lambda| \, \|x\|$; and the triangle inequality $\|x + y\| \le \|x\| + \|y\|$. Expanding the
square gives the frequently used identity
$\|x + y\|^{2} = \|x\|^{2} + 2\, x \cdot y + \|y\|^{2}$.
A unit vector has norm $1$; every $x \neq 0$ is
normalised by $\hat{x} = x / \|x\|$, which keeps the direction and rescales
the length to $1$.
\end{rulebox}

\begin{rulebox}{Angle and Cauchy--Schwarz}
The Cauchy--Schwarz inequality
states that for all $x, y \in \R^{n}$
\[
|x \cdot y| \;\le\; \|x\| \, \|y\|,
\]
with equality iff one vector is a scalar multiple of the other. Hence for
$x, y \neq 0$ the quotient below lies in $[-1, 1]$, and the angle $\theta \in [0, \pi]$ between $x$ and $y$ is defined by
\[
\cos \theta \;=\; \frac{x \cdot y}{\|x\| \, \|y\|}.
\]
Sign reading: $x \cdot y > 0$ means an acute angle, $x \cdot y < 0$ an
obtuse angle, $x \cdot y = 0$ a right angle.
\end{rulebox}

\begin{rulebox}{Orthogonality}
Two vectors are orthogonal, written
$x \perp y$, iff
\[
x \cdot y = 0 .
\]
For orthogonal vectors the identity for $\|x + y\|^{2}$ loses its middle
term and becomes Pythagoras' theorem:
$\|x + y\|^{2} = \|x\|^{2} + \|y\|^{2}$. The zero vector is orthogonal to
every vector.
\end{rulebox}

\begin{rulebox}{Projection}
For $v \neq 0$, the orthogonal projection of $w$ onto $v$ is
\[
\operatorname{proj}_{v}(w) \;=\; \frac{v \cdot w}{\|v\|^{2}} \, v ,
\]
the scalar multiple of $v$ closest to $w$. The residual
$w - \operatorname{proj}_{v}(w)$ is orthogonal to $v$; the projection thus
splits $w$ into a part along $v$ and a part perpendicular to $v$. This
construction reappears in least-squares regression.
\end{rulebox}

\begin{rulebox}{Economic readings}
\begin{itemize}[leftmargin=1.6em,itemsep=0.3ex]
  \item Expenditure: if $p$ is a price vector and $q$ a quantity vector,
        then $p \cdot q = p^{T} q$ is total expenditure; the budget
        constraint with income $m$ reads $p \cdot q \le m$.
  \item Weighted averages: if the weights $w_{i} \ge 0$ sum to $1$, then
        $w \cdot x$ is a weighted average of the entries of $x$ (a convex
        combination), e.g.\ a portfolio return $w^{T} r$ or a course
        grade built from exam scores.
\end{itemize}
\end{rulebox}

\begin{rulebox}{Lines and planes}
\begin{itemize}[leftmargin=1.6em,itemsep=0.3ex]
  \item Line, parametric form: in $\R^{2}$ or
        $\R^{3}$, $x = p + t \, v$ with $t \in \R$, a point $p$ on the
        line, and a direction vector $v \neq 0$.
  \item Line in $\R^{2}$, normal form: $n \cdot x = c$,
        i.e.\ $n_{1} x_{1} + n_{2} x_{2} = c$, with normal vector
        $n \perp v$. Switching: from $v = (v_{1}, v_{2})^{T}$ take
        $n = (-v_{2}, v_{1})^{T}$ and $c = n \cdot p$; from $n$ take
        $v = (-n_{2}, n_{1})^{T}$ and any point $p$ with $n \cdot p = c$.
  \item Plane in $\R^{3}$: parametric form $x = p + s \, v + t \, w$ with
        non-parallel directions $v, w$; normal form $n \cdot x = c$ with
        $n$ orthogonal to both $v$ and $w$ and $c = n \cdot p$.
  \item Distance from a point $q$ to a line in $\R^{2}$ or a plane in
        $\R^{3}$ with normal form $n \cdot x = c$: dividing by the norm
        of the normal,
        \[
        d \;=\; \frac{|n \cdot q - c|}{\|n\|} .
        \]
  \item Intersections lead to small linear systems: for two lines in
        $\R^{2}$, equate the parametric forms (or solve the two normal
        equations), one solution, none (parallel), or infinitely many
        (coinciding). For a line $x = p + t \, v$ and a plane
        $n \cdot x = c$, solve $n \cdot p + t \, (n \cdot v) = c$ for
        $t$; a unique intersection point exists iff $n \cdot v \neq 0$.
\end{itemize}
\end{rulebox}

One bridge remark before drilling: a set of vectors is \emph{linearly
independent} when none of them is a linear combination of the others,
equivalently, when the only combination giving $0$ uses all-zero weights.
That concept, listed alongside the inner product as required review, is
developed and drilled in Chapter~6; the combinations practised here are its
raw material.

\section{Practice Problems}

Use the following vectors throughout:
\[
u = \begin{pmatrix} 1 \\ 2 \\ 2 \end{pmatrix}, \quad
v = \begin{pmatrix} 2 \\ -1 \\ 2 \end{pmatrix}, \quad
w = \begin{pmatrix} 4 \\ 0 \\ 3 \end{pmatrix}, \quad
a = \begin{pmatrix} 3 \\ 4 \end{pmatrix}, \quad
b = \begin{pmatrix} -4 \\ 3 \end{pmatrix}.
\]

\begin{problems}
  \item Compute $u + v$ and $u - w$.
  \item Compute $3 u - 2 v$.
  \item Compute $2 a + b$.
  \item Write $w$ as a linear combination of the standard basis vectors $e_{1}, e_{2}, e_{3}$.
  \item Compute $u \cdot v$, $u \cdot w$, and $v \cdot w$.
  \item Compute $a \cdot b$. What does the result tell you about $a$ and $b$?
  \item Compute $\|u\|$, $\|v\|$, $\|w\|$, and $\|a\|$.
  \item Using $\|\lambda x\| = |\lambda| \, \|x\|$, compute $\|2 w\|$ and $\|-3 u\|$; confirm one of them by direct calculation.
  \item Find the unit vector $\hat{w}$ in the direction of $w$ and check that its norm is $1$.
  \item Find the unit vector in the direction of $c = (6, -8)^{T}$.
  \item Compute $\cos \theta$ for the angle between $u$ and $w$, and give $\theta$ in degrees (two decimals).
  \item Compute the angle between $a$ and $d = (4, 3)^{T}$ (degrees, two decimals).
  \item Compute the angle between $(1, 0)^{T}$ and $(1, 1)^{T}$.
  \item Are $p = (2, 1, -1)^{T}$ and $q = (1, -1, 1)^{T}$ orthogonal? Are $u$ and $v$?
  \item Find $t \in \R$ such that $(t, 2, -1)^{T} \perp (3, t, 4)^{T}$.
  \item Find all $t \in \R$ such that $(t, 3, 2)^{T} \perp (t, t, -5)^{T}$.
  \item Verify the Cauchy--Schwarz inequality numerically for $u$ and $v$. Is it strict here?
  \item Verify the triangle inequality for $a$ and $b$.
  \item Compute $\operatorname{proj}_{v}(w)$ for $v = (1, 2)^{T}$ and $w = (8, 1)^{T}$.
  \item For the data of the previous problem, verify that $w - \operatorname{proj}_{v}(w)$ is orthogonal to $v$.
  \item A consumer buys quantities $q = (10, 4, 6)^{T}$ at prices $p = (3, 2, 5)^{T}$ (in euros). Compute total expenditure $p \cdot q$.
  \item A course grade is the weighted average of scores $x = (90, 70, 80)^{T}$ with weights $\omega = (0.5, 0.3, 0.2)^{T}$. Check that $\omega^{T} x$ is a convex combination of the scores and compute the grade.
\end{problems}

\section{Solutions}

\begin{solutions}
  \item $u + v = (3, 1, 4)^{T}$ and $u - w = (1 - 4,\, 2 - 0,\, 2 - 3)^{T} = (-3, 2, -1)^{T}$.

  \item $3 u - 2 v = (3, 6, 6)^{T} - (4, -2, 4)^{T} = (-1, 8, 2)^{T}$.

  \item $2 a + b = (6, 8)^{T} + (-4, 3)^{T} = (2, 11)^{T}$.

  \item $w = 4 e_{1} + 0 \, e_{2} + 3 e_{3} = 4 e_{1} + 3 e_{3}$.

  \item $u \cdot v = 1 \cdot 2 + 2 \cdot (-1) + 2 \cdot 2 = 2 - 2 + 4 = 4$;
        $u \cdot w = 4 + 0 + 6 = 10$;
        $v \cdot w = 8 + 0 + 6 = 14$.

  \item $a \cdot b = 3 \cdot (-4) + 4 \cdot 3 = -12 + 12 = 0$: the vectors
        are orthogonal, $a \perp b$.

  \item $\|u\| = \sqrt{1 + 4 + 4} = 3$;\;
        $\|v\| = \sqrt{4 + 1 + 4} = 3$;\;
        $\|w\| = \sqrt{16 + 0 + 9} = 5$;\;
        $\|a\| = \sqrt{9 + 16} = 5$.

  \item $\|2 w\| = 2 \cdot 5 = 10$ and $\|-3 u\| = 3 \cdot 3 = 9$. Direct
        check: $2 w = (8, 0, 6)^{T}$, so
        $\|2 w\| = \sqrt{64 + 0 + 36} = \sqrt{100} = 10$.

  \item $\hat{w} = w / \|w\| = \bigl(\tfrac{4}{5}, 0, \tfrac{3}{5}\bigr)^{T}$.
        Check: $\|\hat{w}\|^{2} = \tfrac{16}{25} + 0 + \tfrac{9}{25} = 1$.

  \item $\|c\| = \sqrt{36 + 64} = 10$, so
        $\hat{c} = \bigl(\tfrac{3}{5}, -\tfrac{4}{5}\bigr)^{T}$.

  \item $\cos \theta = \dfrac{u \cdot w}{\|u\| \, \|w\|}
        = \dfrac{10}{3 \cdot 5} = \dfrac{2}{3}$, so
        $\theta = \arccos\bigl(\tfrac{2}{3}\bigr) \approx 48.19^{\circ}$.

  \item $a \cdot d = 12 + 12 = 24$ and $\|a\| = \|d\| = 5$, so
        $\cos \theta = \tfrac{24}{25}$ and
        $\theta = \arccos(0.96) \approx 16.26^{\circ}$.

  \item $\cos \theta = \dfrac{1}{1 \cdot \sqrt{2}} = \dfrac{1}{\sqrt{2}}$,
        so $\theta = 45^{\circ}$.

  \item $p \cdot q = 2 - 1 - 1 = 0$: yes, $p \perp q$. But
        $u \cdot v = 4 \neq 0$, so $u$ and $v$ are \emph{not} orthogonal.

  \item $(t, 2, -1) \cdot (3, t, 4) = 3 t + 2 t - 4 = 5 t - 4 = 0$, so
        $t = \tfrac{4}{5}$.

  \item $(t, 3, 2) \cdot (t, t, -5) = t^{2} + 3 t - 10 = (t + 5)(t - 2)
        = 0$, so $t = 2$ or $t = -5$.

  \item $|u \cdot v| = 4$ and $\|u\| \, \|v\| = 3 \cdot 3 = 9$, and indeed
        $4 \le 9$. The inequality is strict, since $u$ and $v$ are not
        scalar multiples of each other.

  \item $a + b = (-1, 7)^{T}$, so
        $\|a + b\| = \sqrt{1 + 49} = \sqrt{50} = 5 \sqrt{2} \approx 7.07$,
        while $\|a\| + \|b\| = 5 + 5 = 10$. Indeed $7.07 \le 10$.

  \item $v \cdot w = 8 + 2 = 10$ and $\|v\|^{2} = 1 + 4 = 5$, so
        $\operatorname{proj}_{v}(w) = \tfrac{10}{5} \, (1, 2)^{T}
        = (2, 4)^{T}$.

  \item $w - \operatorname{proj}_{v}(w) = (8, 1)^{T} - (2, 4)^{T}
        = (6, -3)^{T}$, and $(6, -3) \cdot (1, 2) = 6 - 6 = 0$.
        Orthogonal, as claimed.

  \item $p \cdot q = 3 \cdot 10 + 2 \cdot 4 + 5 \cdot 6
        = 30 + 8 + 30 = 68$ euros.

  \item The weights are nonnegative and $0.5 + 0.3 + 0.2 = 1$, so
        $\omega^{T} x$ is a convex combination. Grade:
        $0.5 \cdot 90 + 0.3 \cdot 70 + 0.2 \cdot 80 = 45 + 21 + 16 = 82$,
        a value between the smallest score ($70$) and the largest
        ($90$), as every weighted average must be.
\end{solutions}

\section{Additional Problems}

Use the following vectors throughout:
\[
r = \begin{pmatrix} 2 \\ 3 \\ 6 \end{pmatrix}, \quad
s = \begin{pmatrix} 6 \\ 2 \\ 3 \end{pmatrix}, \quad
m = \begin{pmatrix} 5 \\ 12 \end{pmatrix}, \quad
n = \begin{pmatrix} -12 \\ 5 \end{pmatrix}.
\]

\begin{problems}
  \item Compute $r + s$ and $2 r - s$.
  \item Compute $r \cdot s$ and $m \cdot n$.
  \item Compute $\|r\|$, $\|s\|$, and $\|m\|$.
  \item Find the unit vector in the direction of $m$.
  \item Compute $\cos \theta$ for the angle between $r$ and $s$, and give $\theta$ in degrees (one decimal).
  \item Find $t \in \R$ such that $(2, t, 1)^{T} \perp (t, 4, -6)^{T}$.
  \item Find all $t \in \R$ such that $(t, 4, 1)^{T} \perp (t, t, -12)^{T}$.
  \item Compute $\operatorname{proj}_{v}(w)$ for $v = (2, 1)^{T}$ and $w = (3, 4)^{T}$, and verify that the residual $w - \operatorname{proj}_{v}(w)$ is orthogonal to $v$.
  \item Prices are $p = (4, 3, 2)^{T}$ and a consumer wants the bundle $q = (5, 6, 7)^{T}$ with a budget of $50$ euros. Is the bundle affordable? By how much is the budget met or exceeded?
  \item A portfolio puts weights $\omega = (0.25, 0.25, 0.5)^{T}$ on assets with returns $r = (0.04, 0.08, 0.02)^{T}$. Check that the weights form a convex combination and compute the portfolio return $\omega^{T} r$.
  \item Find a parametric form and a normal form of the line $g$ in $\R^{2}$ through the points $P = (1, 2)^{T}$ and $Q = (5, -1)^{T}$.
  \item Compute the distance from the point $q = (4, 6)^{T}$ to the line $g$ of the previous problem.
  \item Decide whether the lines $g\colon x = (1, 1)^{T} + t \, (2, 1)^{T}$ and $h\colon x = (5, 0)^{T} + u \, (-1, 1)^{T}$ intersect, are parallel, or coincide; if they intersect, find the intersection point.
  \item Find the normal form of the plane $E$ in $\R^{3}$ through the point $p = (1, 1, 2)^{T}$ with normal vector $n = (1, 2, 2)^{T}$, and check whether the point $(3, 0, 2)^{T}$ lies on $E$.
  \item Compute the distance from the point $q = (4, 3, 3)^{T}$ to the plane $E$ of the previous problem.
  \item Find the intersection point of the line $x = (2, -1, 1)^{T} + t \, (1, 2, 0)^{T}$ with the plane $E$.
\end{problems}

\section{Additional Solutions}

\begin{solutions}
  \item $r + s = (8, 5, 9)^{T}$ and
        $2 r - s = (4, 6, 12)^{T} - (6, 2, 3)^{T} = (-2, 4, 9)^{T}$.

  \item $r \cdot s = 12 + 6 + 18 = 36$ and
        $m \cdot n = -60 + 60 = 0$ (so $m \perp n$).

  \item $\|r\| = \sqrt{4 + 9 + 36} = \sqrt{49} = 7$;\;
        $\|s\| = \sqrt{36 + 4 + 9} = 7$;\;
        $\|m\| = \sqrt{25 + 144} = \sqrt{169} = 13$.

  \item $\hat{m} = m / \|m\|
        = \bigl(\tfrac{5}{13}, \tfrac{12}{13}\bigr)^{T}$.
        Check: $\tfrac{25}{169} + \tfrac{144}{169} = 1$.

  \item $\cos \theta = \dfrac{36}{7 \cdot 7} = \dfrac{36}{49} \approx
        0.7347$, so $\theta \approx 42.7^{\circ}$.

  \item $(2, t, 1) \cdot (t, 4, -6) = 2 t + 4 t - 6 = 6 t - 6 = 0$, so
        $t = 1$.

  \item $(t, 4, 1) \cdot (t, t, -12) = t^{2} + 4 t - 12 = (t + 6)(t - 2)
        = 0$, so $t = 2$ or $t = -6$.

  \item $v \cdot w = 6 + 4 = 10$ and $\|v\|^{2} = 4 + 1 = 5$, so
        $\operatorname{proj}_{v}(w) = \tfrac{10}{5} \, (2, 1)^{T}
        = (4, 2)^{T}$. Residual: $w - (4, 2)^{T} = (-1, 2)^{T}$, and
        $(-1, 2) \cdot (2, 1) = -2 + 2 = 0$. Orthogonal.

  \item $p \cdot q = 20 + 18 + 14 = 52 > 50$: the bundle is \emph{not}
        affordable; it exceeds the budget by $2$ euros.

  \item The weights are nonnegative and $0.25 + 0.25 + 0.5 = 1$: a convex
        combination. Return:
        $\omega^{T} r = 0.25 \cdot 0.04 + 0.25 \cdot 0.08 + 0.5 \cdot 0.02
        = 0.01 + 0.02 + 0.01 = 0.04 = 4\%$.

  \item Direction: $v = Q - P = (4, -3)^{T}$, so a parametric form is
        \[
        g\colon\; x = \begin{pmatrix} 1 \\ 2 \end{pmatrix}
        + t \begin{pmatrix} 4 \\ -3 \end{pmatrix}, \qquad t \in \R .
        \]
        A normal vector is $n = (3, 4)^{T}$, since
        $n \cdot v = 12 - 12 = 0$, and $c = n \cdot P = 3 + 8 = 11$, so
        the normal form is $3 x_{1} + 4 x_{2} = 11$. Check with $Q$:
        $3 \cdot 5 + 4 \cdot (-1) = 15 - 4 = 11$. Correct.

  \item With $n = (3, 4)^{T}$, $c = 11$, and $\|n\| = \sqrt{9 + 16} = 5$:
        $n \cdot q = 3 \cdot 4 + 4 \cdot 6 = 12 + 24 = 36$, so
        \[
        d = \frac{|n \cdot q - c|}{\|n\|} = \frac{|36 - 11|}{5}
        = \frac{25}{5} = 5 .
        \]

  \item The directions $(2, 1)^{T}$ and $(-1, 1)^{T}$ are not scalar
        multiples of each other ($2 \cdot 1 - 1 \cdot (-1) = 3 \neq 0$),
        so the lines are neither parallel nor coinciding: they intersect
        in exactly one point. Equating components gives the system
        $1 + 2 t = 5 - u$ and $1 + t = u$. Substituting the second
        equation into the first: $1 + 2 t = 5 - (1 + t) = 4 - t$, so
        $3 t = 3$, i.e.\ $t = 1$ and $u = 2$. Intersection point:
        $(1, 1)^{T} + 1 \cdot (2, 1)^{T} = (3, 2)^{T}$. Check on $h$:
        $(5, 0)^{T} + 2 \, (-1, 1)^{T} = (3, 2)^{T}$. Correct.

  \item $c = n \cdot p = 1 \cdot 1 + 2 \cdot 1 + 2 \cdot 2 = 1 + 2 + 4
        = 7$, so
        \[
        E\colon\; x_{1} + 2 x_{2} + 2 x_{3} = 7 .
        \]
        For $(3, 0, 2)^{T}$: $3 + 0 + 4 = 7$, so the point \emph{does}
        lie on $E$.

  \item $\|n\| = \sqrt{1 + 4 + 4} = 3$ and
        $n \cdot q = 4 + 6 + 6 = 16$, so
        \[
        d = \frac{|n \cdot q - c|}{\|n\|} = \frac{|16 - 7|}{3}
        = \frac{9}{3} = 3 .
        \]

  \item Substitute the line into the normal form of $E$: with
        $p = (2, -1, 1)^{T}$ and $v = (1, 2, 0)^{T}$ we get
        $n \cdot p = 2 - 2 + 2 = 2$ and $n \cdot v = 1 + 4 + 0 = 5$, so
        $2 + 5 t = 7$, i.e.\ $t = 1$. Intersection point:
        $(2, -1, 1)^{T} + 1 \cdot (1, 2, 0)^{T} = (3, 1, 1)^{T}$.
        Check: $3 + 2 \cdot 1 + 2 \cdot 1 = 7$. Correct.
\end{solutions}
